\documentclass[a4paper,12pt]{ctexart}
\usepackage{xcolor}
\usepackage[top=1.8cm, bottom=1.8cm, left=1.8cm, right=1.8cm]{geometry}
\usepackage{array}
\usepackage{listings}
\usepackage{hyperref}
\hypersetup{colorlinks=true,linkcolor=blue!70!black}

\lstset{
  basicstyle=\ttfamily\small,
  keywordstyle=\color{blue},
  commentstyle=\color{gray},
  stringstyle=\color{red},
  showstringspaces=false,
  frame=single,
  numbers=left,
  numberstyle=\tiny,
  breaklines=true
}

\title{\textcolor{blue!70!black}{\Huge\textbf{版本控制与协作开发}}}
\author{Git——程序员的"后悔药"和"时光机"}
\date{}

\begin{document}
\maketitle
\thispagestyle{empty}

\section*{写在前面}

你有没有遇到过这种情况：
\begin{itemize}
  \item 改了一堆代码，发现改坏了，但已经回不去了
  \item 想看看一周前的版本，但早就覆盖了
  \item 想和同学一起做项目，但只能一个人写，然后微信传来传去
\end{itemize}

\textbf{Git 就是解决这些问题的工具。} 它是目前全世界程序员都在用的\textbf{版本控制系统}。Git 像一个"时光机"，记录你的每一次修改，可以随时回到任何历史版本。配合 GitHub，你还可以和全世界的人协作。

\section*{第一课：什么是版本控制？}

\subsection*{没有版本控制的时候}

\begin{verbatim}
我的论文_初稿.doc
我的论文_改了一版.doc
我的论文_改了两版.doc
我的论文_最终版.doc
我的论文_最终版2.doc
我的论文_打死也不改了.doc
\end{verbatim}

——是不是很眼熟？这就是没有版本控制的惨状。

\subsection*{有了 Git 之后}

\begin{verbatim}
Git 记录：
commit 1: "初稿完成"
commit 2: "修改了第三章"
commit 3: "加了参考文献"
commit 4: "改了几个错别字"
\end{verbatim}

每个 \texttt{commit}（提交）就是一次快照。你随时可以跳回任何一个 commit 查看当时的文件状态。

\subsection*{核心概念}

\begin{center}
\begin{tabular}{|c|l|}
\hline
\textbf{概念} & \textbf{简单理解} \\ \hline
仓库（Repository） & 你的项目文件夹 + Git 的版本记录 \\ \hline
提交（Commit） & 一次快照——"到此一游"的记录 \\ \hline
分支（Branch） & 平行宇宙——在不同分支上做不同的事 \\ \hline
合并（Merge） & 把两个平行宇宙合并到一起 \\ \hline
远程仓库（Remote） & 存在 GitHub 上的备份，方便协作 \\ \hline
\end{tabular}
\end{center}

\section*{第二课：安装与配置}

\subsection*{安装 Git}

\begin{itemize}
  \item \textbf{Windows:} 下载 \url{https://git-scm.com/download/win}，一路默认安装
  \item \textbf{Mac:} 打开终端，输入 \texttt{git}，系统会提示安装
\end{itemize}

安装完成后，打开命令行验证：

\begin{lstlisting}[language=bash]
git --version
\end{lstlisting}

\subsection*{配置身份}

Git 需要知道你是谁（用于记录谁做了修改）：

\begin{lstlisting}[language=bash]
git config --global user.name "你的名字"
git config --global user.email "你的邮箱@example.com"
\end{lstlisting}

\section*{第三课：Git 基本工作流}

\subsection*{初始化仓库}

\begin{lstlisting}[language=bash]
cd 你的项目文件夹
git init
\end{lstlisting}

这会在文件夹里创建一个隐藏的 \texttt{.git} 目录——Git 的所有信息都存在这里。

\subsection*{三区概念}

\begin{verbatim}
工作区（Working Directory）    暂存区（Staging Area）    仓库（Repository）
   你的文件                     准备提交的文件           已提交的历史
        │                            │                       │
        │   git add <file>           │   git commit          │
        └───────────────────────────→└──────────────────────→┘
\end{verbatim}

\subsection*{日常工作流}

\begin{lstlisting}[language=bash]
# 1. 查看当前状态（重要！你应该经常运行）
git status

# 2. 添加文件到暂存区
git add 文件名          # 添加特定文件
git add .               # 添加所有修改过的文件

# 3. 提交（创建快照）
git commit -m "提交说明：改了什么东西"

# 4. 查看历史
git log                 # 查看提交历史
git log --oneline       # 简版历史
\end{lstlisting}

\subsection*{练习：创建你的第一个 Git 仓库}

\begin{enumerate}
  \item 新建一个文件夹 \texttt{my-first-repo}
  \item 用 \texttt{git init} 初始化
  \item 创建一个 \texttt{hello.py} 文件，写入 \texttt{print("Hello Git!")}
  \item \texttt{git add hello.py}
  \item \texttt{git commit -m "第一个提交"}
  \item 修改 \texttt{hello.py}，加一行 \texttt{print("这是第二次修改")}
  \item \texttt{git status} 看看有什么变化
  \item \texttt{git add .} 然后 \texttt{git commit -m "加了第二行"}
  \item \texttt{git log --oneline} 看历史
\end{enumerate}

\section*{第四课：时光机——回到过去}

\subsection*{查看修改}

\begin{lstlisting}[language=bash]
# 查看尚未暂存的修改
git diff

# 查看已经暂存的修改
git diff --staged
\end{lstlisting}

\subsection*{撤销修改}

\begin{lstlisting}[language=bash]
# 放弃工作区的修改（回到最近一次 commit 的状态）
git checkout -- 文件名

# 从暂存区撤回到工作区
git reset HEAD 文件名

# 撤回最近一次 commit
git reset --soft HEAD~1     # 保留工作区的修改
git reset --hard HEAD~1     # 彻底回到上一个版本（危险！）
\end{lstlisting}

\subsection*{回到特定版本}

\begin{lstlisting}[language=bash]
# 先用 git log 找到你想回去的 commit id
git log --oneline
# 输出：4a3b2c1 修改了标题
#       f7e6d5c 第一个提交

# 回到特定版本（--hard 会丢掉之后的修改）
git reset --hard 4a3b2c1

# 或者只是看看当时的样子
git checkout 4a3b2c1
# 看完后回到最新版：
git checkout main
\end{lstlisting}

\section*{第五课：分支——平行宇宙}

分支是 Git 最强大的功能。你可以在一个分支上开发新功能，同时在另一个分支上修 bug。

\begin{verbatim}
main ────●────●────●──────────●─────
              \              /
feature       ●────●────●──
\end{verbatim}

\subsection*{分支操作}

\begin{lstlisting}[language=bash]
# 查看当前分支（* 表示当前所在分支）
git branch

# 创建新分支
git branch feature-login

# 切换到新分支
git checkout feature-login

# 创建并切换一步到位
git checkout -b feature-login

# 合并分支（先切回 main）
git checkout main
git merge feature-login

# 删除分支
git branch -d feature-login
\end{lstlisting}

\subsection*{练习：分支工作流}

\begin{enumerate}
  \item 在 main 分支上创建一个 \texttt{README.md}，写项目简介
  \item 创建并切换到 \texttt{feature-add} 分支
  \item 在 feature-add 分支上添加一个 \texttt{calculator.py}
  \item 切回 main 分支——看看 \texttt{calculator.py} 还在吗？（不在！）
  \item 把 feature-add 合并到 main
  \item 现在 \texttt{calculator.py} 出现了
\end{enumerate}

\subsection*{合并冲突}

如果两个分支修改了同一个文件的同一行，Git 不知道该用哪个版本——这就叫\textbf{冲突}。

\begin{lstlisting}[language=bash]
# 合并时出现冲突，Git 会说：CONFLICT
# 打开文件你会看到：
<<<<<<< HEAD
这是 main 分支的版本
=======
这是 feature 分支的版本
>>>>>>> feature

# 手动修改成你想要的样子，然后：
git add 文件名
git commit -m "解决了冲突"
\end{lstlisting}

\section*{第六课：GitHub——远程协作}

\subsection*{什么是 GitHub？}

GitHub 是一个存放 Git 仓库的网站。你可以把代码备份到上面，也可以和他人协作。

\begin{itemize}
  \item \url{https://github.com/} —— 注册一个账号
\end{itemize}

\subsection*{SSH 密钥配置}

为了让你的电脑和 GitHub 安全通信，需要配置 SSH 密钥：

\begin{lstlisting}[language=bash]
# 生成密钥（一路回车即可）
ssh-keygen -t rsa -b 4096 -C "你的邮箱"

# 查看公钥
cat ~/.ssh/id_rsa.pub
\end{lstlisting}

把输出的内容复制到 GitHub → Settings → SSH and GPG keys → New SSH key。

\subsection*{推送代码到 GitHub}

\begin{lstlisting}[language=bash]
# 在 GitHub 上新建一个仓库
# 然后把本地仓库推送到远程

# 添加远程仓库
git remote add origin git@github.com:你的用户名/仓库名.git

# 推送到远程
git push -u origin main

# 以后只需要
git push
\end{lstlisting}

\subsection*{从 GitHub 拉取代码}

\begin{lstlisting}[language=bash]
# 克隆别人的仓库
git clone git@github.com:用户名/仓库名.git

# 拉取最新修改
git pull
\end{lstlisting}

\section*{第七课：Pull Request——协作的标准流程}

在实际团队协作中，你不会直接把代码推送到 main 分支。标准流程是：

\subsection*{协作工作流}

\begin{enumerate}
  \item 从 main 创建一个新分支：\texttt{git checkout -b my-feature}
  \item 在新分支上写代码、commit
  \item 推送到远程：\texttt{git push origin my-feature}
  \item 在 GitHub 上创建一个 \textbf{Pull Request}（PR）
  \item 团队成员 review 你的代码
  \item 通过后，合并到 main
  \item 删除 feature 分支
\end{enumerate}

\subsection*{模拟协作练习}

\begin{enumerate}
  \item 两人一组（或自己用两个 GitHub 账号）
  \item A 创建仓库，添加 B 为 collaborator
  \item B 克隆仓库，创建分支，修改代码，推送，创建 PR
  \item A review 代码，合并 PR
  \item B 切回 main，\texttt{git pull} 获取最新代码
\end{enumerate}

\section*{第八课：.gitignore——不要什么都存}

有些文件不应该被 Git 管理：
\begin{itemize}
  \item 编译生成的文件（\texttt{.exe}, \texttt{.o}）
  \item 依赖包（\texttt{node\_modules}, \texttt{\_\_pycache\_\_}）
  \item 密码和密钥
  \item 大文件
\end{itemize}

在项目根目录创建 \texttt{.gitignore} 文件：

\begin{lstlisting}
# Python
__pycache__/
*.pyc
*.pyo
.env

# C++
*.exe
*.o
*.obj

# IDE
.vscode/
.idea/

# OS
.DS_Store
Thumbs.db
\end{lstlisting}

\section*{第九课：Git 实用技巧}

\subsection*{查看谁写了这段代码}

\begin{lstlisting}[language=bash]
git blame 文件名
# 显示每一行是谁在什么时候写的
\end{lstlisting}

\subsection*{暂存当前工作}

如果你写到一半，需要切换到其他分支处理紧急问题：

\begin{lstlisting}[language=bash]
git stash                # 把工作区暂存起来
git stash pop            # 恢复暂存的工作
\end{lstlisting}

\subsection*{修改最后一次 commit}

\begin{lstlisting}[language=bash]
# 如果 commit message 写错了，或者漏了一个文件
git commit --amend -m "新的提交说明"
\end{lstlisting}

\subsection*{查看美观的历史}

\begin{lstlisting}[language=bash]
git log --graph --oneline --all
# 以图形方式显示所有分支的历史
\end{lstlisting}

\subsection*{常用命令速查}

\begin{center}
\begin{tabular}{|c|c|}
\hline
\textbf{命令} & \textbf{作用} \\ \hline
\texttt{git init} & 初始化仓库 \\ \hline
\texttt{git clone <url>} & 克隆远程仓库 \\ \hline
\texttt{git add <file>} & 添加到暂存区 \\ \hline
\texttt{git commit -m "msg"} & 提交 \\ \hline
\texttt{git status} & 查看状态 \\ \hline
\texttt{git log} & 查看历史 \\ \hline
\texttt{git diff} & 查看修改 \\ \hline
\texttt{git branch} & 查看分支 \\ \hline
\texttt{git checkout -b <name>} & 创建并切换分支 \\ \hline
\texttt{git merge <branch>} & 合并分支 \\ \hline
\texttt{git push} & 推送到远程 \\ \hline
\texttt{git pull} & 从远程拉取 \\ \hline
\texttt{git stash} & 暂存当前工作 \\ \hline
\texttt{git reset --hard <id>} & 回到历史版本 \\ \hline
\end{tabular}
\end{center}

\section*{总结：Git 的核心思想}

\begin{verbatim}
本地                             远程
┌─────────────┐                ┌─────────────┐
│ 工作区       │   git push    │  GitHub      │
│ 暂存区       │ ───────────→  │  (备份+协作)  │
│ 本地仓库     │   git pull    │              │
│  (历史版本)  │ ←───────────  │              │
└─────────────┘                └─────────────┘
\end{verbatim}

\textbf{你学到了：}
\begin{itemize}
  \item Git 的三区概念（工作区→暂存区→仓库）
  \item 基本工作流（add → commit → push）
  \item 时光回溯（reset, checkout）
  \item 分支与合并
  \item GitHub 远程协作
  \item Pull Request 工作流
  \item .gitignore、stash、blame 等技巧
\end{itemize}

Git 是程序员的\textbf{必修课}。从现在开始，你写的每一个项目都应该用 Git 管理。一开始可能会不习惯，但用几周后你就会发现——没有 Git 的日子简直无法想象。

\section*{综合练习与答案}

\subsection*{基础题}

\begin{enumerate}
  \item 创建一个 Git 仓库，进行 5 次提交，然后用 \texttt{git log --oneline --graph} 查看历史。
  \item 创建一个分支 \texttt{test-branch}，在上面修改一个文件，然后合并回 main。
  \item 模拟一个合并冲突：在两个不同的分支上修改同一文件的同一行，然后合并。
  \item 用 \texttt{git stash} 暂存工作，切换到另一个分支做修改，再切回来恢复暂存。
  \item 创建一个 \texttt{.gitignore} 文件，忽略 \texttt{.log} 文件和 \texttt{temp/} 目录。
\end{enumerate}

\subsection*{参考答案}

\textbf{练习 3（模拟合并冲突）：}
\begin{verbatim}
# 在 main 分支上
echo "这是 main 分支的内容" > conflict.txt
git add conflict.txt
git commit -m "main 版本"

# 创建并切换到 feature 分支
git checkout -b feature
echo "这是 feature 分支的内容" > conflict.txt
git add conflict.txt
git commit -m "feature 版本"

# 切回 main 合并
git checkout main
git merge feature
# 输出：CONFLICT 冲突！
# 手动编辑 conflict.txt 保留想要的内容
git add conflict.txt
git commit -m "解决冲突"
\end{verbatim}

\section*{扩展项目}

\subsection*{项目一：用 Git 管理你的博客项目}

把 \texttt{web-dev/} 模块的博客项目用 Git 管理起来。

\textbf{要求：}
\begin{enumerate}
  \item 初始化仓库，每次完成一个功能就 commit
  \item 用分支开发新功能（如"文章分类"功能开一个分支）
  \item 推送到 GitHub
  \item 写一个完善的 README.md
  \item 邀请一个朋友做 Code Review
\end{enumerate}

\subsection*{项目二：开源贡献初体验}

在 GitHub 上找一个适合初学者的开源项目（搜索 "good first issue" 标签），尝试：
\begin{enumerate}
  \item fork 项目
  \item 修复一个文档 typo 或简单的 bug
  \item 提交 Pull Request
  \item 等待 maintainer 合并
\end{enumerate}

\section*{本模块与其他模块的关联}

\begin{itemize}
  \item \textbf{Web 开发} ↔ 本模块：\texttt{web-dev/} 模块的博客项目是你练习 Git 的最佳实操对象。用分支管理功能开发、用 GitHub 管理代码、用 Pull Request 做代码审查——这是真实开发团队的标准工作流。
  \item 本模块 → \textbf{综合项目实战}：capstone 模块要求你独立完成一个完整项目，而 Git 是项目管理的基础工具。你在本模块学会的版本控制技能会贯穿 capstone 的整个开发过程。
  \item \textbf{Python 基础} → 本模块：本模块所有命令行操作在计算思维模块的"命令⾏基础"部分有铺垫。
\end{itemize}

\end{document}
